
\section*{\underline{EXPOS\'E XIII}\\[3pt] \underline{\underline{LES TH\'EOR\`EMES DE FINITUDE POUR LE FONCTEUR DE PICARD}}}
\begin{center}
\underline{par S. KLEIMAN}
\end{center}

\subsection*{1. Les (b)-faisceaux}
Soient @@M0@@ un corps alg\'ebriquement clos et @@M1@@ un @@M2@@-sch\'ema projectif, muni d'un module inversible ample @@M3@@.

\textbf{D\'efinition 1.1 (Mumford).} \emph{Soit @@M4@@. Un faisceau coh\'erent @@M5@@ sur @@M6@@ est dit \emph{@@M7@@-r\'egulier} si, en chaque point du support de @@M8@@, @@M9@@ est engendr\'e par ses sections globales et si, pour tout @@M10@@, on a}
@@M11@@

\textbf{Lemme 1.2.} \emph{Soit @@M12@@ une suite exacte de faisceaux coh\'erents sur @@M13@@}\SGAAuthorFootnote{*}{Il est sous-entendu, ici comme dans la suite dans des situations analogues, que @@M14@@ est d\'efini par tensorisation par une section donn\'ee de @@M15@@.}\emph{. Si @@M16@@ est @@M17@@-r\'egulier, alors @@M18@@ l'est.}

En effet, on a la suite exacte correspondante :
@@M19@@

\textbf{Proposition 1.3.} \emph{Si un faisceau @@M20@@ sur @@M21@@ est @@M22@@-r\'egulier, alors, pour tout @@M23@@, on a}~:
\begin{enumerate}[label=(\roman*)]
\item \emph{@@M24@@ est @@M25@@-r\'egulier.}
\item \emph{@@M26@@ est surjectif.}
\item \emph{@@M27@@ est engendr\'e par @@M28@@.}
\end{enumerate}

